\documentclass[a4paper, 14pt]{article}
\usepackage{a4wide}
\usepackage[utf8]{inputenc}
\usepackage[russian]{babel}
\usepackage[dvips]{graphicx, color}
\usepackage{epstopdf}
\usepackage{amsmath}
\usepackage{amsfonts}
\usepackage{fancyref}
\usepackage{indentfirst}
\usepackage{extsizes}
\usepackage{amssymb}
\usepackage{amsthm}
\usepackage{mathtools}
\usepackage{tikz}
\usepackage{pgfplots}
\usepackage{multirow} % улучшенное форматирование таблиц
\usepackage{ulem} % подчеркивания
\usepackage{geometry}
\usepackage{breqn}
\geometry{left=3cm}
\geometry{right=1.5cm}
\geometry{top=2.4cm}
\geometry{bottom=2.4cm}
\DeclareMathOperator*{\argmax}{\arg\!\max}

\begin{document}


\linespread{1.3} % полуторный интервал
%\renewcommand{\rmdefault}{ftm} % Times New Roman
\frenchspacing
\thispagestyle{empty}


\begin{center}
\ \vspace{-3cm}

\includegraphics[width=0.5\textwidth]{msu.eps}\\

{\scshape Московский государственный университет имени М.~В.~Ломоносова}\\
Факультет вычислительной математики и кибернетики\\
Кафедра системного программирования

\vfill

{\LARGE Курсовая работа}

\vspace{1cm}

{\Huge\bfseries <<Определение языка сообщений социальной сети Twitter>>} \\

\end{center}

\vspace{1cm}

\begin{flushright}
  \large
  \textit{Выполнил студент 327 группы}\\
  М.~А.~Кольцов

  \vspace{5mm}

  \textit{Научный руководитель}\\
  В.~Д.~Майоров
\end{flushright}

\vfill

\begin{center}
Москва, 2014
\end{center}

\pagebreak
\tableofcontents
\pagebreak

\begin{center}Аннотация\end{center}

\hangindent=3.5cm \hangafter=0 \noindent В данной работе рассматривается задача определения языка сообщений социальной сети Twitter: с помощью имеющейся выборки твитов нужно построить классификатор, который будет для новых текстов выдавать предполагаемый язык. Сравниваются три самостоятельных системы для решения этой задачи с тремя авторскими реализациями, основанными на актуальных статьях. Показывается конкурентоспособность имеющихся решений, а также пригодность к решению задачи одного из предложенных.

\newtheorem{definition} {Определение}
\newtheorem{option} {Свойство}
\newtheorem{theorem} {Теорема}


\section{Введение}
	        В настоящее время человечество имеет доступ к огромному запасу знаний, накопленному за тысячи лет. Немалая часть этих знаний представляется
	        в виде текстов на различных языках. В связи с этим активно разрабатываются методы, предназначенные для автоматического извлечения
	        и преобразования информации, данной в символьном представлении. 
	        Возникло научное направление <<обработка естественных языков>>. Одной из его фундаментальных задач является определение языка текста. 


	        Стандартным подходом к этой задаче является применение машинного обучения. А именно, если у нас есть база из сотен документов 
	        на нескольких языках, то можно <<предсказать>> язык поступившего на рассмотрение документа, сравнив его с имеющимися. 
	        В общем случае, нужно по имеющимся данным построить так называемую модель, а затем все действия с текстами проводить в терминах 
	        этой модели. 
	        Классическим примером является метод, описанный в 1994 году: каждому документу сопоставим <<профиль документа>> --- упорядоченный по
	        числу вхождений список n-грамм --- последовательностей длины n подряд идущих символов. <<Профиль языка>> --- это совокупность профилей
	        документов, которые имеются на этом языке. Теперь, если нужно для какого-то документа определить язык, то последовательность
	        действий такова:
	        \begin{enumerate}
	        		\item Составляется профиль документа
	        		\item Происходит сравнение полученного профиля с имеющимися профилями языков
	        		\item Тот язык, чей профиль наиболее похож на профиль документа, объявляется результатом
	        \end{enumerate}
	          
	        Вышеописанный метод плохо работает для коротких сообщений. В то же время, поток информации в виде коротких неформальных сообщений нельзя игнорировать ---
	        в социальной сети Twitter среднее количество сообщений в день составляет более пятидесяти восьми миллионов\footnote{По данным http://www.statisticbrain.com/twitter-statistics/ на 1 января 2014 года}, а длина каждого ограничена 140 символами. 
	        Такой формат текстов обусловил появление новых алгоритмов. В данной работе рассматриваются современные методы 
	        решения задачи определения языка, а также предлагается улучшение для одного из них. Проводится тестирование, показывающее
	        превосходство по качеству распознавания языка полученного результата над существующими.
  
	        
\section{Постановка задачи}
		\subsection{Формальное описание задачи автоматического определения языка}
		Пусть $L$ - множество меток, сопоставленных естественным языкам.
		По заданному тренировочному корпусу $$ T = \{(msg_{1}, lang_{1}), (msg_{2}, lang_{2}), \ldots, (msg_{N}, lang_{N}) \} $$
		($msg_{i}, i \in \overline{1..N},$ - текст на естественном языке, $lang_{i} \in L$ - метка этого языка) нужно построить классификатор,
		который произвольному входному сообщению $new\_msg$ на языке $some\_lang$ сопоставит метку $l \in L$, соответствующую этому языку,
		 или же сообщит, что язык текста невозможно достоверно распознать.
		\subsection{Цели и задачи курсовой работы}
		В данной работе в качестве текстов выступают так называемые <<твиты>> --- сообщения из социальной сети Twitter\footnote{https://twitter.com/}, 
		а множество $L$ соответствует 18 языкам, которые можно разделить на несколько групп по типу алфавита:
		\begin{itemize}
			\item Основанные на кириллице: болгарский, чувашский, русский, татарский, украинский
			\item Арабские: арабский, персидский (фарси), урду
			\item Латинские: нидерландский, французский, английский, немецкий, итальянский, испанский, турецкий
			\item Деванагари: хинди, маратхи, непальский
		\end{itemize}
		Цели работы:
		\begin{enumerate}
			\item Исследовать современные решения задачи автоматического определения языка коротких сообщений
			\item Провести совместное сравнительное тестирование некоторых методов решения задачи и выяснить, действительно ли они показывают
			конкурентноспособное качество классификации
			\item Исследовать зависимость качества классификации от различных факторов: степени нормализации сообщений, мощности обучающей 
			выборки
			\item Исследовать возможность улучшения какого-либо алгоритма решения задачи автоматического определения языка коротких сообщений	
		\end{enumerate}
		В рамках выполнения цели работы необходимо решить следующие задачи:
		\begin{enumerate}
			\item Собрать обучающий корпус, состоящий не менее чем из пятисот твитов каждого языка
			\item Реализовать один или несколько методов решения задачи
			\item Выбрать адекватную имеющейся задаче меру качества классификации
		\end{enumerate}

\section{Обзор существующих решений}
		\subsection{Подход, основанный на подсчёте частоты n-грамм}
		Для каждого языка составляется список из $K$ самых часто встречающихся n-грамм. Для $new\_msg$ составляется аналогичный список, для него
		определяется наиболее <<похожий>> из списков языков
		и на основании этого делается вывод о языке сообщения. 
		
		<<Похожесть>> двух упорядоченных списков характеризуется метрикой \textit{out-of-place}: 
		пусть в одном списке n-грамма $x$ стоит на позиции $i$, а в другом \nolinebreak --- $j$, тогда к расстоянию между списками добавляется $|i - j|$.
		
		
		Cavnar \& Trenkle выдвинули гипотезу, что $ \approx 300$ самых часто используемых n-грамм сильно зависят
		от языка, и экспериментально подтвердили её. Также они предлагают использовать длину n-грамм вплоть до 5.
		
		Этим подходом пользуются многие программные продукты, о которых речь пойдёт
		в разделе 4.
		
		\noindent Плюсы и минусы:
		
		
		$+$ Небольшой размер модели
		
		$+$ Высокая скорость обучения и классификации
		
		$-$ Качество классификации сильно зависит от длины текста
		
		\subsection{Prediction by partial matching}
		Составляется кодирование, основываясь на количествах вхождений в документы цепочек из не более чем $n$ символов (более коротким кодовым словам соответствуют более часто встречающиеся последовательности). Затем для $new\_msg$ вычисляется длина кода при помощи полученного кодирования и выбирается наиболее
		компактно кодирующий язык.
		
		Такой подход предлагается, например, в работе \cite{ppm}.
		
		\noindent Плюсы и минусы:
		
		$+$ 	Один параметр --- $n$ --- для настройки метода
		
		$+$ Показывает высокое качество классификации
		
		$-$ Нуждается в больших объёмах обучающих данных
		
		\subsection{Подходы, использующие алгоритмы машинного обучения}
		Из документов выделяются признаки, характеризующие их, и каждому документу сопоставляется вектор признаков. 
		Таким образом, получается задача классификации в общем виде. Поскольку это более широкая задача, она была
		лучше изучена и есть множество известных подходов для её решения. Один из примеров --- 
		метод опорных векторов.
		 
		Благодаря гибкости в выборе признаков, метод является перспективным по отношению к определению языка твитов. Так, например, в работах \cite{ppm} \cite{lrev}
		используются такие признаки: 
		\begin{itemize}
			\item количество вхождений n-грамм
			\item имя пользователя в Twitter и его местоположение
			\item язык интерфейса
			\item наиболее вероятный язык в предыдущих $k$ твитах этого пользователя
			\item наиболее вероятный язык сущностей, на который ссылается пользователь в сообщении: текстов
		других пользователей, ссылок, тэгов.
		\end{itemize}		
		
		\noindent Плюсы и минусы:
		
		$+$ Разработанный математический аппарат
		
		$+$ Возможность адаптации для выбранной предметной области
		
		$-$ Нужно тщательно подбирать признаки и алгоритм, дабы получить хорошее качество классификации
		
		\subsection{Language identification graph-based approach (LIGA)}
		По каждому документу обучающей выборки строится граф $(V, E)$:
		 каждой вершине $v \in V$ графа соответствует пара $(trigram, count)$, где $trigram$ --- это одна
		из триграмм, встретившихся в тексте, а $count$ - количество её вхождений. Будем обозначать за $v_{tr}$ соответствующую вершине $v$ триграмму
		($v_{tr} \neq u_{tr} \ \forall u \in V/ \{v\}$), за $v_{cnt}$ - количество её вхождений. 
		Ориентированное ребро $(u, v) \in E$ соответствует тому, что в тексте $u_{tr}$ следует непосредственно
		перед $v_{tr}$, а $(u, v)_{cnt}$ --- количество таких вхождений.
		
		Далее, выполняется объединение полученных графов для всех документов одного языка $l \in L$: если у двух графов $(V', E')$ и $(V'', E'')$ есть вершины
		$v' \in V'$ и $v'' \in V''$ такие, что $v'_{tr} = v''_{tr}$, то в результирующем графе $(V_{l}, E_{l})$ должна быть вершина 
		$v \in V_{l}$ такая, что 
		$$v_{tr} = v'_{tr} = v''_{tr}$$ $$v_{cnt} = v'_{cnt} + v''_{cnt}$$
		Аналогично выполняется объединение рёбер.
		
		Последний шаг в построении модели: объединить графы $(V_{l}, E_{l})$, полученные $\forall l \in L$, в один. 
		Для этого каждой вершине $v \in V_{l}$
		(ребру $(u, v) \in E_{l}$) ставится дополнительно в соответствие метка языка $v_{lang}$ \ ($(u, v)_{lang}$). 
		Готовая модель представляет собой граф
		($\mathcal{V}, \mathcal{E}$), который получается следующим образом:
		$$ \mathcal{V} = \bigcup_{l \in L} V_{l}$$
		$$ \mathcal{E} = \bigcup_{l \in L} E_{l}$$
		
		Классификация происходит декомпозицией сообщения на триграммы, вычислением функции схожести сообщения с языком и
		выбором языка с максимальным значением этой функции. Формально, если
		$t_{1}, t_{2}, \ldots, t_{n}$ - триграммы сообщения $new\_msg$ в порядке их вхождения, то результатом будет
		$$ \argmax_{l} score_{l} ,$$
		где вектор \textit{score} длины $|L|$ определяется как:
		$$ score_{l} = \sum_{i=1}^{n} get\_v(t_{i}, l) + \sum_{i=1}^{N-1} get\_e(t_{i}, t_{i+1}, l)$$
		\[
 		get\_v(tr, l) =
 		  \begin{dcases}
  		   v_{cnt} \over \sum_{w \in \mathcal{V}} w_{cnt} & 
  		   \begin{split} 
					 if \ \exists v \in \mathcal{V}: v_{tr}=tr \\ v_{lang}=l	
  		   	\end{split} \\
  		   0 & \quad otherwise
  		 \end{dcases}
		\]		
		
		\[
 		get\_e(tr_1, tr_2, l) =
 		  \begin{dcases}
  		   (u, v)_{cnt} \over \sum_{(w, q) \in \mathcal{E}} (w, q)_{cnt} & 
  		    \begin{split} 
					 if \ \exists (u, v) \in \mathcal{E}:  u_{tr}=tr_1 \\ v_{tr}=tr_2 \\ (u, v)_{lang}=l
  		   	\end{split} \\
  		   0 & \quad otherwise
  		 \end{dcases}
		\]	
		
		Такая модель позволяет зафиксировать одновременно частотность триграмм и их положение, при этом все триграммы
		из документов участвуют в классификации. Например, неправильное написание какого-то слова, будучи упущенным в n-граммном подходе (поскольку он
		учитывает лишь наиболее часто встречающиеся сочетания), добавит определённости при определении языка методом LIGA.		
		
		\noindent Плюсы и минусы:
		
		$+$ Высокая скорость обучения и классификации
		
		$+$ Максимально использует информацию из обучающей выборки
		
		$+$ Возможность дообучения <<на лету>>
		
		$+$ Возможность визуализации модели
		
		$-$ Большой объём модели и, как следствие, относительно высокие требования к памяти
			
\section{Исследование и построение решения задачи}
		\subsection{Сравниваемые системы}
			\subsubsection{TextCat}
			TextCat\footnote{http://odur.let.rug.nl/vannoord/TextCat/} --- реализация подхода Canvar \& Trenkle \cite{canvar}, основанного на n-граммах.
			Поскольку этот метод был одним из первых, практически все появляющиеся системы сравниваются именно
			с ним. 
			Оригинальные модели языков достаточно устарели, но, к счастью, у программы есть возможность обучения.
			\subsubsection{Google CLD}
			Google compact language detector\footnote{https://code.google.com/p/cld2/} --- программный продукт, встроенный в браузер Chromium, 
			предназначенный для определения языка просматриваемой страницы. 
			Использует 4-граммы и умеет считывать метаинформацию
			о веб-страницах. Впрочем, для коротких текстов он тоже используется, например, в сервисе Google Translate. Система поставляется обученной
			на огромном корпусе текстов на более чем 100 языках.
			\subsubsection{Langid.py}
			Программа langid.py\footnote{https://github.com/saffsd/langid.py} за авторством Lui \& Baldwin \cite{langid}, выпущенная в 2011 году, позиционируется
			как быстрое и точное решение задачи определения языка текста. Продукт поддерживает 97 языков.
			\subsubsection{LIGA}
			Реализация метода LIGA, выполненная в рамках курсовой работы в соответствии со статьёй \cite{liga}. Подход подробно описан в секции 3.4.
			\subsubsection{Предлагаемое улучшение LIGA}
			Я предлагаю улучшенную версию алгоритма LIGA. Основное отличие заключается в изменении функций get\_v и get\_e:
			\[
 			get\_v\_v2(tr, l) =
 			  \begin{dcases}
  			   {(\log{|L| \over |\{r \in L | \exists v \in \mathcal{V}: v_{tr}=tr, v_{lang}=r\}|} + 1) \times  v_{cnt}} \over 
  			   \sum_{w \in \mathcal{V},\ w_{lang}=l} w_{cnt} & 
  			   \begin{split} 
						 \exists v \in \mathcal{V}: v_{tr}=tr \\ v_{lang}=l	
  			   	\end{split} \\
  			   0 & \quad otherwise
  			 \end{dcases}
			\]		
		
			\[
 			get\_e\_v2(tr_1, tr_2, l) =
 			  \begin{dcases}
  			   {(\log{|L| \over edges} + 1) \times (u, v)_{cnt}} 
  			   \over \sum_{(w, q) \in \mathcal{E},\ (w, q)_{lang}=l} (w, q)_{cnt} & 
  			    \begin{split} 
						 \exists (u, v) \in \mathcal{E}:  u_{tr}=tr_1 \\ v_{tr}=tr_2 \\ (u, v)_{lang}=l
  			   	\end{split} \\
  			   0 & \quad otherwise
  			 \end{dcases}
			\]	
			$$ edges = |\{r \in L \ | \ \exists (u, v) \in \mathcal{E}: u_{tr}=tr_1, v_{tr}=tr_2, (u, v)_{lang}=r\}| $$
			Мотивация к такому изменению следующая. Гораздо полезнее знать, насколько специфична данная триграмма данному языку, чем то, сколько раз
			 она встречалась в обучающей выборке вообще. Для этого в знаменателе общая сумма очков всех вершин графа (аналогично для рёбр) заменена 
			на сумму очков всех вершин подграфа $(V_{l}, E_{l})$. Также вводится логарифмический коэффициент, который характеризует <<уникальность>> данной триграммы в обучающей выборке.
			
			После таких изменений появляется шанс, что фраза \textit{"Привiт, груша"} будет корректно распознана как украинская, поскольку триграммы
			\textit{вiт} и \textit{ивi} будут иметь больший вес, чем остальные.
			\subsubsection{LogR}
			Реализация метода LogR, выполненная в рамках курсовой работы в соответствии со статьёй \cite{ppm}. Общий подход описан в секции 3.3. Данная версия использует логистическую регрессию из пакета LibLinear \cite{fan}, а также следующие признаки:
			\begin{itemize}
				\item Логарифм количества вхождений каждой уни-, би-, три- и квадграммы.
				\item Имя оставившего твит пользователя Twitter, его местоположение и отображаемое имя, а также их префиксы
				\item Индикатор: написано ли имя пользователя латиницей
				\item Встреченные в твите хэштеги
					\footnote{Хэштег - это последовательность символов без пробелов, начинающаяся с '$\#$', например: \textit{\#HarryPotter}, \textit{\#russia2014}}
				\item Встреченные в твите упоминания
					\footnote{Упоминание - это отображаемое имя пользователя, предварённое символом '@', например: \textit{@KremlinRussia}}
				\item Доменное имя 1 и 2 уровней, а также протокол для всех ссылок, встреченных в сообщении (при этом, если ссылка была сокращённой,
				например bit.ly/something, то происходит http-запрос для определения конечной цели)
			\end{itemize}
		\subsection{Обучающий корпус}   
			Твиты для обучающего корпуса собирались из следующих мест:
			\begin{enumerate}
				\item \textbf{Из статей по теме распознавания языка в Twitter.} В работе \cite{liga} тексты сообщений оказались в открытом доступе, около 9 тысяч на западноевропейских языках. Также авторы статьи \cite{ppm} предоставили собранный корпус, в котором около 6 тысяч твитов, содержащих помимо текста ещё и метаинформацию о пользователях. Огромный набор текстов на русском языке прилагается к статье \cite{julia} - более двухсот тысяч твитов.
				\item \textbf{Извлечённые с помощью TwitterAPI по идентификационному номеру.} Согласно политике Twitter, в открытом доступе могут находиться не тексты сообщений, а лишь их идентификационные номера. Вместе с работой \cite{lrev} распространяется архив, содержащий номера твитов, используемых в ней при тестировании. 
				\item \textbf{С сайта Indigenous Tweets.}\footnote{http://indigenoustweets.com/} Данная веб-страница содержит список малопредставленных в Twitter языков, а также перечисление всех известных пользователей с указанием количества сообщений, которое они оставили на конкретном языке. Были выделены пользователи, пишущие в основном на татарском и чувашском языках, а их сообщения были выкачаны с помощью TwitterAPI. Таким образом было получено 6 тысяч твитов.
				\item \textbf{С помощью собственного аккаунта.} В распоряжении оказался аккаунт турецкого пользователя, подписанного на новостные рассылки на турецком языке. Было извлечено 800 твитов на турецком. Такое малое количество обусловлено ограничениями TwitterAPI.
			\end{enumerate}		
			Итоговое распределение по языкам (после нормализации по схеме 2) показано в таблице 1.	  
			 
			\begin{center}
			\begin{table}
			\begin{tabular*}{\textwidth}{|@{\extracolsep{\fill} }l  r|}
				\hline
				Язык  & Количество твитов \\
				\hline
				Арабский & 1171 \\
				Английский & 2234 \\
				Болгарский & 1880 \\
				Испанский & 2350 \\
				Итальянский & 1538 \\
				Маратхи & 1154 \\
				Немецкий & 2208 \\
				Непальский & 1678 \\
				Нидерландский & 2032 \\
				Персидский (фарси) & 2361 \\
				Русский & 196836 \\
				Татарский & 3248 \\
				Турецкий & 781 \\
				Украинский & 627 \\
				Урду & 1073 \\
				Французский & 2239 \\
				Хинди & 1209 \\
				Чувашский & 2584 \\
				\hline
				Всего & 227203 \\
				\hline
			\end{tabular*}
			\caption{Распределение по языкам твитов обучающей выборки.}
			\end{table}
			\end{center}		
			
		\subsection{Описание сценариев нормализации}	
		Сообщения в социальных сетях часто содержат ошибки, выражения эмоций, неправильное употребление слов и знаков препинания. Чтобы классификатор
		был устойчивым к такого рода явлением, нужно \textit{нормализовывать}	 каждое сообщение --- то есть приводить его к какому-то стандартному виду.
		Для того, чтобы проследить зависимость качества классификации от степени нормализации текстов, были использованы следующие сценария нормализации:
		\begin{enumerate}
			\item Удаляются отметки о ретвитах\footnote{Ретвит --- твит, отправитель которого не является автором, при этом обычно в тексте идёт имя автора после букв RT, например, RT @KremlinRussia}, ссылки, упоминания, хэштеги, символы пунктуации заменены на пробел, цифры удалены, последовательно идущие пробельные символы заменены на один пробел
			\item Всё то же, что и в пункте один, плюс: все последовательности из трёх и более подряд идущих одинаковых символов заменены на два символа (например, слово \textit{оооооооу} превращается таким образом в \textit{ооу}, что помогает сгладить последствия проявления эмоций), все слова из 2 и менее букв удалены, текст переведён в нижний регистр.
		\end{enumerate}
		В обоих случаях сообщения, в которых после нормализации не осталось никаких символов, не включались в итоговую выборку.
		
		\subsection{Выбор меры качества классификации}
		В данной работе в качестве основной меры качества используется \textit{полнота}, то есть доля правильно определённых сообщений на фиксированном языке
		от общего количества твитов на этом языке.
		
		Выбор обусловлен тем, что классификатор LIGA, ввиду большого числа обучающих данных на русском языке,
		выдаёт почти всем кириллическим текстам метку <<русский>>. В то же время, на несколько десятков, например, болгарских сообщений он выдаёт 
		<<болгарский>>. Если посчитать \textit{точность} определения болгарского языка --- отношение числа правильно определённых болгарских твитов 
		к числу твитов, определённых как болгарские --- то она будет близка к 100\%. Точность распознавание русского языка при этом падает.
		Но, так как русских сообщений на два порядка больше, то падает она примерно на 1\%. Таким образом, высокие показатели \textit{точности} на деле
		скрывают неудовлетворительные результаты классификации.
		
		Тем не менее, были вычислены все общепринятые меры качества мультиклассовой классификации: то
       
\subsection{Результаты тестирования}
		Тестирование проводилось при помощи одной из разновидностей метода \textit{кросс-валидации}. 
		Обучающая выборка делилась на две части: тренировочную и тестовую. Те классификаторы, у которых есть
		возможность обучиться на имеющихся данных, обучались на тренировочной части. Затем собирались результаты
		предсказаний языка на тестовой части выборки, вычислялись точность, полнота и F-мера. Такое разбиение
		происходило 10 раз для каждого классификатора, а полученные результаты усреднялись.
		
		Реализация LogR обучалась на два порядка дольше других классификаторов (во многом из-за того, что делала много http-запросов
		на раскрытие ссылок), а при уменьшении числа твитов в обучающей выборке показывала неудовлетворительный результат, близкий к случайной
		выдаче языка, поэтому не вошла в итоговые таблицы и графики.
		
		Прочерки в таблицах означают, что классификатор ни разу не выдал в качестве ответа такой язык.
		
		\subsubsection{Зависимость полноты классификации от количества данных, даваемых на обучение}
		Таблица 2 показывает, что вышеуказанной зависимости не наблюдается.
		
		\subsubsection{Зависимость полноты классификации от степени нормализации}
		На таблице 3 видно, что в полнота, в целом, падает при использовании сценария 2 для нормализации. Впрочем, 
		стоит заметить, что алгоритм LIGAv2 преодолел барьер в 80\% для украинского языка.

\begin{center}
\begin{table}
\begin{tabular*}{1.07\textwidth}{| l| *{5}{|c} |}
\hline 
Язык & Google CLD2 & langid.py & LIGAv2 & LIGA & TextCat \\
\hline
\hline
Болг. & 81/81/81/82 & 77/77/77/77 & 94/94/94/94 & 0/0/0/0 & 88/88/88/87 \\
Рус. & 91/91/91/91 & 86/86/86/86 & 90/90/90/90 & 99/99/99/99 & 87/87/87/87 \\
Татар. & -/-/-/- & -/-/-/- & 99/99/99/99 & 00/00/00/00 & 99/99/99/99 \\
Укр. & 80/82/81/81 & 88/88/88/88 & 76/78/77/79 & -/-/-/- & 88/88/88/89 \\
Чуваш. & -/-/-/- & -/-/-/- & 98/99/98/99 & 00/00/00/00 & 96/96/96/96 \\
\hline
\hline
Араб. & 84/82/84/84 & 90/90/90/90 & 97/97/98/98 & 89/89/90/89 & 95/95/94/95 \\
Перс. & 87/87/87/88 & 89/89/89/88 & 94/94/94/94 & 95/95/95/95 & 88/88/88/88 \\
Урду & 89/89/89/89 & 90/90/89/90 & 97/96/96/96 & 81/81/82/82 & 92/92/92/92 \\
\hline
\hline
Англ. & 97/97/97/97 & 96/96/96/96 & 87/88/88/87 & 92/92/92/92 & 93/93/93/93 \\
Испан. & 86/86/86/85 & 93/93/93/93 & 97/96/97/96 & 97/96/96/96 & 91/91/91/91 \\
Итал. & 92/92/92/92 & 98/98/98/98 & 98/98/98/98 & 92/92/92/92 & 96/96/97/96 \\
Нем. & 94/94/94/94 & 96/96/96/97 & 97/97/97/97 & 95/95/95/95 & 94/93/93/93 \\
Нидер. & 90/90/91/90 & 96/96/96/96 & 98/98/98/98 & 96/96/96/96 & 94/94/94/94 \\
Турец. & 88/88/89/88 & 89/90/89/89 & 94/94/94/93 & 14/14/14/15 & 89/89/89/88 \\
Франц. & 91/92/91/92 & 96/96/96/96 & 96/96/96/96 & 97/98/97/98 & 93/93/93/93 \\
\hline
\hline
Марат. & 91/91/91/91 & 80/81/80/82 & 92/91/91/91 & 73/73/73/74 & 81/82/80/80 \\
Непал. & 88/88/88/88 & 71/71/70/70 & 92/93/92/92 & 93/93/94/93 & 81/82/83/83 \\
Хинди & 97/97/97/97 & 83/83/84/85 & 95/95/95/95 & 85/86/86/86 & 82/83/81/82 \\
\hline
\end{tabular*}
\caption{Полнота классификации при обучении на 50/60/70/80\% от общего числа твитов на данном языке.}
\end{table}
\end{center}

\begin{center}
\begin{table}
\begin{tabular*}{0.933\textwidth}{| l| *{5}{c} |}
\hline 
Язык & Google CLD2 & langid.py & LIGAv2 & LIGA & TextCat \\
\hline
\hline
Болгарский & 81/84 & 77/70 & 94/92 & 00/00 & 88/87 \\
Русский & 91/93 & 86/86 & 90/94 & 99/99 & 87/86 \\
Татарский & -/- & -/- & 99/99 & 00/00 & 99/99 \\
Украинский & 81/83 & 88/87 & 77/84 & -/- & 88/83 \\
Чувашский & -/- & -/- & 98/90 & 00/00 & 96/95 \\

\hline
\hline
Арабский & 84/78 & 90/92 & 98/98 & 90/90 & 94/93 \\
Персидский & 87/80 & 89/87 & 94/92 & 95/95 & 88/84 \\
Урду & 89/85 & 89/79 & 96/93 & 82/63 & 92/88 \\

\hline
\hline
Английский & 97/95 & 96/92 & 88/85 & 92/88 & 93/90 \\
Испанский & 86/81 & 93/84 & 97/94 & 96/93 & 91/82 \\
Итальянский & 92/90 & 98/95 & 98/98 & 92/84 & 97/95 \\
Немецкий & 94/93 & 96/96 & 97/97 & 95/97 & 93/91 \\
Нидерландский & 91/90 & 96/93 & 98/97 & 96/91 & 94/90 \\
Турецкий & 89/91 & 89/88 & 94/94 & 14/27 & 89/86 \\
Французский & 91/88 & 96/92 & 96/95 & 97/95 & 93/97 \\
\hline
\hline
Маратхи & 91/90 & 80/83 & 91/94 & 73/79 & 80/75 \\
Непальский & 88/88 & 70/70 & 92/93 & 94/94 & 83/80 \\
Хинди & 97/93 & 84/72 & 95/85 & 86/64 & 81/75 \\

\hline
\end{tabular*}
\caption{Полнота классификации при использовании первого/второго сценария нормализации.}
\end{table}
\end{center}

		\subsubsection{Примеры ошибок}
		Выявлены следующие причины возникновения ошибок (в скобках после сообщений указан сначала предполагаемый язык, а затем предсказанный):
		\begin{enumerate}
			\item Похожесть языков: \textit{lang leve ikea} (nl $\to$ en), 
				\textit{state department condemns concerted campaign intimidate international journalists cairo} (en $\to$ fr).
				Некоторые языки с родственным алфавитом похожи. Не имеющий знаний об арабских языках
				 человеку сложно отличить один арабский язык от другого. Хоть применение
				машинного обучения и помогает в этом, но некоторые конструкции всё же не поддаются различию.
			\item Наличие имён собственных на другом языке:
				\textit{moyes ponders toffees selection everton boss david moyes will have decisions make both ends the pitch} (en $\to$ fr).
				Названия групп, фирм, мест --- всё это накладывает свой отпечаток на сообщение, и в этом случае классификаторам
				легко ошибиться даже на достаточно длинных сообщениях.
			\item Фразы на другом языке:
				\textit{харесах видеоклип watch mario balotelli failed backheel} (bg $\to$ en). В таких твитах сложно определить, какой
				язык является первичным. Например, некоторые твиты на языке урду представляют из себя записи вида примерно 
				<<learn new word --- self --- in urdu: ...>>. Однозначный вывод сделать автоматически очень сложно. Возможно, в этом
				помог бы анализ метаинформации о пользователе.
			\item Неверная разметка: 
				\textit{встиг підстригтися уже наступний раз записався} (ru $\to$ uk). Некоторые тексты, особенно среди двухсоттысячного набора 
				русских твитов, на самом деле не написаны на предполагаемом языке. К счастью, их не очень много и они мало влияют на оценку
				качества, тем более что все классификаторы на них <<ошибаются>>
			\item Малая длина:
				\textit{кис}, \textit{прекратите}, \textit{lesles}, \textit{haa dus}, \textit{lool grave}. Тексты этой категории сложно отнести
				к конкетному языку даже живому человеку. При ужесточении нормализации таких коротких твитов становится больше
		\end{enumerate}
	
		\subsubsection{Выводы}
		тут некие выводы 
		
		

\section{Описание практической части}
		тут я пишу о том, как запускать мою программу

\pagebreak
	
\section{Заключение}
		В рамках данной работы предполагалось изучить актуальные методы решения задачи автоматического определения языка сообщений социальной сети Twitter, провести сравнительное тестирование некоторых, исследовать возможность улучшения качества классификации. 
		
		Поставленные цели были достигнуты и были получены следующие результаты:
		\begin{enumerate}
			\item Реализовано 2 подхода к решению этой задачи на языке программирования \textit{Python}, а также модернизированная версия одного из них
			\item Собрана корпус из более чем двухсот тысяч сообщений
			\item Написан комплекс скриптов для проведения совместного тестирования полученных решений с программными продуктами, успешно используемыми в области обработки естественых языков
			\item Проведено тестирование, сделаны выводы на основе результатов
		\end{enumerate}
		Была исследована зависимость качества классификации от степени нормализации сообщений, 
		рассмотрены примеры ошибочно определённых твитов, сделаны предположения по поводу причин ухудшения качества работы классификаторов.
		
		Тестирование показало, что улучшенная версия метода LIGA не уступает, а зачастую и превосходит по качеству классификации все остальные методы.
		
		

  \begin{thebibliography}{99}
    \bibitem{canvar} Cavnar W. B. et al. N-gram-based text categorization //Ann Arbor MI. – 1994. – Т. 48113. – №. 2. – С. 161-175.
    \bibitem{liga} Tromp E., Pechenizkiy M. Graph-based n-gram language identification on short texts //Proc. 20th Machine Learning conference of Belgium and The Netherlands. – 2011. – С. 27-34.
    \bibitem{ppm} Bergsma S. et al. Language identification for creating language-specific Twitter collections //Proceedings of the Second Workshop on Language in Social Media. – Association for Computational Linguistics, 2012. – С. 65-74.
    \bibitem{lrev} Carter S., Weerkamp W., Tsagkias M. Microblog language identification: Overcoming the limitations of short, unedited and idiomatic text //Language Resources and Evaluation. – 2013. – Т. 47. – №. 1. – С. 195-215.
    \bibitem{julia} Ю.В. Рубцова. Метод построения и анализа корпуса коротких текстов для задачи классификации отзывов // Электронные библиотеки: перспективные методы и технологии, электронные коллекции: Труды XV Всероссийской научной конференции RCDL’2013, Ярославль, Россия, 14-17 октября 2013 г. – Ярославль: ЯрГУ, 2013. –С. 269-275.
    \bibitem{fan} Fan R. E. et al. LIBLINEAR: A library for large linear classification //The Journal of Machine Learning Research. – 2008. – Т. 9. – С. 1871-1874.
    \bibitem{langid} Baldwin T., Lui M. Language identification: The long and the short of the matter //Human Language Technologies: The 2010 Annual Conference of the North American Chapter of the Association for Computational Linguistics. – Association for Computational Linguistics, 2010. – С. 229-237.
  \end{thebibliography}
\end{document}
